\PassOptionsToPackage{noend}{algorithm2e}
\documentclass[pmlr]{jmlr}% new name PMLR (Proceedings of Machine Learning Research)

 % The following packages will be automatically loaded:
 % amsmath, amssymb, natbib, graphicx, url, algorithm2e

 %\usepackage{rotating}% for sideways figures and tables
\usepackage{longtable}% for long tables

 % The booktabs package is used by this sample document
 % (it provides \toprule, \midrule and \bottomrule).
 % Remove the next line if you don't require it.
\usepackage{booktabs}
 % The siunitx package is used by this sample document
 % to align numbers in a column by their decimal point.
 % Remove the next line if you don't require it.
\usepackage[load-configurations=version-1]{siunitx} % newer version
 %\usepackage{siunitx}

 % The following command is just for this sample document:
\newcommand{\cs}[1]{\texttt{\char`\\#1}}

 % Define an unnumbered theorem just for this sample document:
\theorembodyfont{\upshape}
\theoremheaderfont{\scshape}
\theorempostheader{:}
\theoremsep{\newline}
\newtheorem*{note}{Note}

 % change the arguments, as appropriate, in the following:
\jmlrvolume{1}
\jmlryear{2026}
\jmlrworkshop{Workshop Title}


% Han's packages
\usepackage{tikz}
\usetikzlibrary{arrows.meta, positioning, matrix, tikzmark}
\usepackage{subcaption}
\newcommand{\refs}{%
    \textcolor{red}{\textbf{[NEED REF]}}%
    }
\usepackage[table]{xcolor}
\usepackage{todonotes}




\title{Association Transformations in Multi-tier Structures}

%\title[Short Title]{Full Title of Article\titlebreak This Title Has
%A Line Break\titletag{\thanks{sample footnote}}}
%
 % Use \Name{Author Name} to specify the name.

 % Spaces are used to separate forenames from the surname so that
 % the surnames can be picked up for the page header and copyright footer.
 
 % If the surname contains spaces, enclose the surname
 % in braces, e.g. \Name{John {Smith Jones}} similarly
 % if the name has a "von" part, e.g \Name{Jane {de Winter}}.
 % If the first letter in the forenames is a diacritic
 % enclose the diacritic in braces, e.g. \Name{{\'E}louise Smith}

 % *** Make sure there's no spurious space before \nametag ***

 % Two authors with the same address
%  \author{\Name{Author Name1\nametag{\thanks{with a note}}} \Email{abc@sample.com}\and
%   \Name{Author Name2} \Email{xyz@sample.com}\\
%   \addr Address}

 % Three or more authors with the same address:
 % \author{\Name{Author Name1} \Email{an1@sample.com}\\
 %  \Name{Author Name2} \Email{an2@sample.com}\\
 %  \Name{Author Name3} \Email{an3@sample.com}\\
 %  \Name{Author Name4} \Email{an4@sample.com}\\
 %  \Name{Author Name5} \Email{an5@sample.com}\\
 %  \Name{Author Name6} \Email{an6@sample.com}\\
 %  \Name{Author Name7} \Email{an7@sample.com}\\
 %  \Name{Author Name8} \Email{an8@sample.com}\\
 %  \Name{Author Name9} \Email{an9@sample.com}\\
 %  \Name{Author Name10} \Email{an10@sample.com}\\
 %  \Name{Author Name11} \Email{an11@sample.com}\\
 %  \Name{Author Name12} \Email{an12@sample.com}\\
 %  \Name{Author Name13} \Email{an13@sample.com}\\
 %  \Name{Author Name14} \Email{an14@sample.com}\\
 %  \addr Address}


 % Authors with different addresses:
 % \author{\Name{Author Name1} \Email{abc@sample.com}\\
 % \addr Address 1
 % \AND
 % \Name{Author Name2} \Email{xyz@sample.com}\\
 % \addr Address 2
 %}

%\editor{Editor's name}
 % \editors{List of editors' names}

\begin{document}

\maketitle

%\begin{abstract}
%This is the abstract for this article.
%\end{abstract}
%\begin{keywords}
%List of keywords
%\end{keywords}


\section{Introduction}
This paper presents a method for learning transformations over 
multi-tier representations, with a particular focus on modifications 
of cross-tier relations (i.e., associations between tiers). The 
multi-tiered structures are widely attested in linguistics and are 
known as \textit{Autosegmental Representations}\citep{CG1984autosegmental,goldsmith1976autosegmental,hayes1980metrical,sagey1986representation}, which
are graph structures of temporally-ordered 
tiers and cross-tier association lines. This work focuses on a 
specific but important type of transformation where the elements
on each tier remain unchanged from input to output, while the 
association lines are modified subject to non-crossing constraints. 
We formalize these transformations as graph model structures \citep{strother2019diss} 
and use a bottom-up factor inference method 
\citep{chandleeBufia2019,rawski2021structure} to infer which 
association lines are preserved, deleted, or newly introduced from 
observed input-output pairs. 
 
\begin{figure}[ht!]
	\centering
	
	\begin{tikzpicture}[
		node/.style={inner sep=1.5pt, minimum size=16pt},
		assoc/.style={thick},
		scale=0.9
		]
		
		% ================= INPUT =================
		
		% top tier
		\node[node] (a1) at (0,1.5) {$a_1$};
		\node[node] (a2) at (1.2,1.5) {$a_2$};
		\node[node] (a3) at (2.4,1.5) {$a_3$};
		
		% bottom tier
		\node[node] (b1) at (0,0) {$b_1$};
		\node[node] (b2) at (1.2,0) {$b_2$};
		\node[node] (b3) at (2.4,0) {$b_3$};
		
		% associations
		\draw[assoc] (a1) -- (b1);
		\draw[assoc] (a2) -- (b2);
		\draw[assoc] (a3) -- (b3);
		
		% subcaption
		\node at (1.2,-0.8) {Input};
		
		% ================= ARROW =================
		
		\draw[->, thick] (3.5,0.75) -- (4.7,0.75);
		
		% ================= OUTPUT =================
		
		% top tier
		\node[node] (c1) at (5.8,1.5) {$a_1$};
		\node[node] (c2) at (7.0,1.5) {$a_2$};
		\node[node] (c3) at (8.2,1.5) {$a_3$};
		
		% bottom tier
		\node[node] (d2) at (7.0,0) {$b_2$};
		\node[node] (d3) at (8.2,0) {$b_3$};
		
		% associations
		\draw[assoc] (c1) -- (d2);
		\draw[assoc] (c2) -- (d3);
		\draw[assoc] (c3) -- (d3);
		
		% subcaption
		\node at (7.0,-0.8) {Output};
		
	\end{tikzpicture}
	
	\caption{
		Learning association transformations between tiers.
	}
\end{figure}

Such multi-tiered structures have been widely used in linguistic 
theory, such as tone-syllable associations in autosegmental 
phonology \citep{goldsmith1976autosegmental}
prosodic or syllable structure \citep{hayes1980metrical,
	sagey1986representation}. 
The core idea behind autosegmental graphs is that linguistic 
information has multiple parallel information organized across 
temporally ordered tiers. The input-output 
mappings are in fact more naturally characterized by rearranging 
association lines across tiers rather than changing elements 
themselves. A large body of theoretical work on African tonal 
processes has demonstrated the advantages of autosegmental graphs 
\citep{CG1984autosegmental, odden2013introducing}.  
Computationally, using multi-tiered structures has been shown to 
reduce the complexity of tonotactic patterns \citep{jardine2016computationally,Chandlee-jardine2021input}.
Previous work has modeled these structures using finite-state 
approaches \citep{wiebe1992modelling,bird1994one,yli2013finite,
	yli2015three}. These approaches typically convert structured 
representations into strings or synchronized tapes to leverage 
finite-state machinery, or rely on hard-coded mechanisms to 
represent different types of associations.

In contrast, we propose an approach that learns association 
transformations between tiers directly over autosegmental 
graphs using model theory and formal logic. As a preliminary 
study, we begin with an idealized assumption: tonal processes 
involve no changes to segmental content, but only the 
rearrangement of association lines between the underlying 
representation (UR) and the surface representation (SR). 
In addition, we assume that only a single process constitutes 
the learning target. 

We demonstrate that, by adopting autosegmental graphs within a 
model-theoretic framework of formal transductions, it is 
possible to construct a partial order over structures and 
apply a model-theoretic learner such as the Bottom-Up Factor 
Inference Algorithm \citep{chandleeBufia2019,rawski2021structure}. 
The learner outputs a formula that preserves 
unchanged associations, deletes associations that must be delinked, and 
introduces newly added associations. Learning proceeds through iterative 
hypothesis updates while traversing a partially ordered space of logically possible structures.

Furthermore, transformations that appear non-deterministic when modeled 
over strings can instead be expressed deterministically within this 
relational framework. To our knowledge, the problem of learning local 
transformations of association relations in multi-tier structures has not previously been addressed in this form.


\section{Preliminaries}
\subsection{Autosegmental graphs}

ASGs represent the idea that different aspects of linguistic structure
are organized on distinct tiers while progressing simultaneously over time.
Each tier encodes a particular type of information,
such as tone (pitch information), segmental (tongue shape), syllabic
(metrical) structure. 
Note that there may be multiple tiers in ASGs, provided that the 
elements on different tiers are disjoint. In most phonological 
applications, ASGs typically involve two or three tiers. In this 
work, we assume a two-tiered structure, though the definition 
remains the same in the general case. In addition, ASGs are subject 
to the No Crossing Constraint (NCC), which enforces temporal 
alignment between associations.

\begin{definition}[Autosegmental Graph]
	An autosegmental graph (ASG) is a labeled $\mathcal{T}$-partite graph
	whose vertices are partitioned into $\mathcal{T}$ disjoint sets, with 
	each tier linearly-ordered. Equivalently, it can be defined
	model-theoretically as a relational structure
	$\mathcal{A} = \langle D, \mathcal{T}, \alpha, \lhd \rangle$
	where $D$ is a finite domain of elements; $\mathcal{T} = \{T_1,\dots,T_n\}$ is a 
	collection of tiers such that
	$T_1 \cup T_2 \cup \dots \cup T_n = D$ and $T_i \cap T_j = \emptyset$
	for all $i \neq j$. 
	$\lhd$ is a binary successor relation defined between nodes within tiers;
	$\alpha \subseteq D \times D$ is a set of non-crossing association relations across tiers. 
\end{definition}

\begin{definition}[No Crossing Constraint (NCC)]
	An ASG needs to obey No Crossing Constraint. Formally,
	For all $(x,y), (x',y') \in \alpha$,
	\(
	x \lhd x' \Rightarrow (y \lhd y') \vee (y = y').
	\) and vice versa.
\end{definition}


\begin{figure}[ht]
	\centering
	
	% ================= AUTSEG =================
	\begin{minipage}[c]{0.20\textwidth}
		\centering
		
		\begin{tikzpicture}
			\node (v1) at (0,0) {$\sigma$};
			\node (v2) at (1,0) {$\sigma$};
			\node (t1) at (0,1) {L};
			\node (t2) at (1,1) {H};
			
			\draw (v1)--(t2);			
			\draw (v1)--(t1);
			\draw (v2)--(t2);
		\end{tikzpicture}
		
	\end{minipage}
	\hspace{0.03\textwidth}
	% ================= GRAPH =================
	\begin{minipage}[c]{0.20\textwidth}
		\centering
		
		\begin{tikzpicture}[
			node/.style={draw, circle, inner sep=2pt, minimum size=18pt},
			assoc/.style={thick},
			scale=0.8
			]
			
			\node[node] (t1) at (0,2) {1\textsubscript{L}};
			\node[node] (t2) at (2,2) {2\textsubscript{H}};
			
			\node[node] (s1) at (0,0) {3\textsubscript{$\sigma$}};
			\node[node] (s2) at (2,0) {4\textsubscript{$\sigma$}};
			
			\draw (t1) -- node[midway, left] {$\alpha$} (s1);
			\draw (t2) -- node[midway, right] {$\alpha$} (s2);
			\draw (t2) -- node[midway, right] {$\alpha$} (s1);	
			
			\draw (t1) -- node[midway, above] {$\lhd$} (t2);
			\draw (s1) -- node[midway, above] {$\lhd$} (s2);
			
		\end{tikzpicture}
		
	\end{minipage}
	\hspace{0.03\textwidth}
	% ================= MODEL =================
	\begin{minipage}[c]{0.40\textwidth}
		\centering
		\footnotesize
		
		\[
		\begin{aligned}
			D &= \{1,2,3,4\} \\
			T &= \{1,2\}\quad L = \{1\}\quad H = \{2\} \\
			\sigma &= \{3,4\} \\
			\lhd &= \{
			\langle 1,2 \rangle,
			\langle 3,4 \rangle
			\} \\
			\alpha &= \{
			\langle 1,3 \rangle,
			\langle 2,4 \rangle,
			\langle 2,3 \rangle
			\}
		\end{aligned}
		\]
		
	\end{minipage}

	\caption{
		Three representations of a two-tier autosegmental graph.
	}
	\label{fig:armodel}	
\end{figure}

Figure~\ref{fig:armodel} shows a tone--syllable structure and its
corresponding model. In this example, the first syllable is associated
with both a Low (L) tone and a following High (H) tone, while the second
syllable is associated with the H tone (often
represented as LH.H in linguistic literature).
The model consists of four labelled nodes, numbered \(1,2,3,4\), each
representing a linguistic unit with its corresponding properties. There
are two tiers: the tone tier \(T\), containing the L and H tones, and the
syllable tier \(\sigma\), containing the two syllables. The successor
relation, denoted by \(\lhd\), represents ordering within a single tier.
In this example, the successor relation is
$\lhd = \{\langle 1,2 \rangle,\ \langle 3,4 \rangle\}$.
The association relation, denoted by \(\alpha\), represents cross-tier
associations:
$
\alpha =
\{\langle 1,3 \rangle,\ \langle 2,3 \rangle,\ \langle 2,4 \rangle\}.
$.
Together, the tuple
$\langle D,T,\sigma,L,H,\lhd,\alpha\rangle$
constitutes the \textit{signature} of this ASG model.

The size of an autosegmental graph is defined as the sum
of both node and association counts, excluding successor relations,
since successor count is implicitly determined by the count
of nodes on each tier. For example, the size of the ASG in 
Figure \ref{fig:armodel} is 7.

\begin{definition}[Size of ASG]
Given $\mathcal{A}=\langle D, \mathcal{T},\alpha,\lhd\rangle$, 
the size of \(\mathcal{A}\), denoted by \(|\mathcal{A}|\), 
is defined as $\mathcal{A} = |D| + |\alpha|$.
\end{definition}

For convenience, we also define the notion of neighbors in an ASG.
Two nodes are considered neighbors if they are connected by a
successor relation.
\begin{definition}[Neighbor]
Given $\mathcal{A}=\langle D, \mathcal{T},\alpha,\lhd\rangle$, 
a node \(n\) is a neighbor of a node \(m\) iff
$m \lhd n \lor n \lhd m$.
\end{definition}

\subsection{Matrix Representation}
In addition to representing model-theoretic relational structures 
with graphs, one can use \textit{adjacency matrices} which use binary values to
denote whether there is an edge between two nodes. Two nodes
are considered \textit{adjacent} if there is an edge connecting
them. For example, the ASG shown in Figure~\ref{fig:armodel}
can be equivalently represented as in Table \ref{tab:matx}.
\begin{figure}[ht!]
\centering

\begin{definition}[Adjacency matrices]
Given an ASG
\(
\mathcal{A}=\langle D,\mathcal{T},\alpha,\lhd\rangle
\)
with \(t\) tiers, the association relation \(\alpha\) can be represented
by \(t-1\) adjacency matrices, each corresponding to the
association relation between two adjacent tiers
\(T_i,T_j \in \mathcal{T}\).
The adjacency matrix
\(M_{\alpha}^{(i,j)}\)
is defined such that
\[
M_{\alpha}^{(i,j)}[k,\ell] =
\begin{cases}
1 & \text{if } (x_k,y_\ell)\in\alpha, \\
0 & \text{otherwise},
\end{cases}
\qquad
\begin{array}{c|cc}
     & L & H \\
\hline
\sigma_1 & 1 & 1 \\
\sigma_2 & 0 & 1 \\
\end{array}
\]
where \(x_k \in T_i\) and \(y_\ell \in T_j\).
\end{definition}

\caption{
Adjacency matrix representation of the ASG in
Figure~\ref{fig:armodel}.
}
\label{tab:matx}

\end{figure}

The motivation for adopting matrices is to establish a framework
corresponding to string-based learning. Concepts such as substrings,
subfactors, and superfactors are well developed and widely used in
grammar inference \refs, but have not been extensively explored for
ASGs. To make these concepts applicable to ASGs, we follow the notions of
\textit{restriction} and \textit{containment} proposed in
\citet{chandleeBufia2019}, adapting them to graph-based
representations through the notions of \textit{submatrices} and \textit{matrix
containment.}

\begin{definition}[Submatrix]
\label{def.submatrix}
\textit{
Let matrix \(A\) be an \(m \times n\) matrix.
A matrix \(B\) is a \emph{submatrix} of \(A\) if it is formed by taking
\(k\ (1 \leq k \leq m)\) contiguous rows and
\(v\ (1 \leq v \leq n)\) contiguous columns from \(A\).
}
\end{definition}

\begin{definition}[Containment]
\textit{
An \(m \times n\) matrix \(A=[a_{ij}]\) is contained
in a matrix \(B\) (written \(A \sqsubseteq B\))
if there exists an \(m \times n\) submatrix of \(B\),
denoted \(B'=[b'_{ij}]\),
such that for all
\(i \in [1,m]\) and \(j \in [1,n]\),
\(
b'_{ij} \geq a_{ij}.
\)
}
\end{definition}

Now that we can determine whether a matrix \(A\) is a submatrix of
\(B\), we can also determine whether \(A\) is contained in \(B\).
If so, \(A\) is a subfactor of \(B\), while \(B\) is a
superfactor of \(A\). Furthermore, if the size difference between
a superfactor and its subfactor is exactly one, then the superfactor
is called the \textit{next superfactor} of that subfactor.

Containment is important because it provides a partial ordering
over the space of ASG factors, which enables the learner to infer
entailment relations. The idea is simple: if  structure forms
the foundation of the learning algorithm, which determines whether
one structure is contained in another. Importantly, all factors 
that contain a forbidden structure, i.e., its superfactors, are 
also forbidden, while all factors contained in a well-formed 
structure, i.e., its subfactors, are likewise well formed.
%
%
%\begin{figure}
%\centering
%\begin{tikzpicture}[
%	every node/.style={inner sep=1pt, font=\small\itshape},
%	edge/.style={},
%	]
%	\matrix (m) [matrix of nodes,
%	row sep=6mm,
%	column sep=10mm,
%	nodes={anchor=center},
%	]{
%%   \node (aaa) {aaa}; &\node (llb1) {...};&   \node (llb2) {...}; & \node (llb3) {...};&\node (llb4) {...};&\node (llb5) {...};&    \node (bbb) {bbb}; \\
%   &\node (aa) {aa}; & \node (ab) {ab}; && \node (ba) {ba}; & \node (bb) {bb}; \\
%	 && \node (a) {a}; & & \node (b) {b}; &  \\
%	 &&& \node (z) {$\epsilon$}; && \\};
%
%	\draw[edge] (z) -- (a);
%
%	\draw[edge] (z) -- (b);
%	
%	\draw[edge] (a) -- (aa) ;
%	\draw[edge] (a) -- (ab);
%	\draw[edge] (a) -- (ba);
%	
%	\draw[edge] (b) -- (bb) ;
%	\draw[edge] (b) -- (ab);
%	\draw[edge] (b) -- (ba);
%	\draw[edge, dotted] (aa) -- ++(0,1.5em);
%%	\draw[edge, dotted] (aa) -- ++(1,1.5em);
%	\draw[edge, dotted] (ab) -- ++(0,1.5em);
%%	\draw[edge, dotted] (ab) -- ++(1,1.5em);	
%%	\draw[edge, dotted] (bb) -- ++(-1,1.5em);	
%	\draw[edge, dotted] (ba) -- ++(0,1.5em);
%	\draw[edge, dotted] (bb) -- ++(0,1.5em);
%	\end{tikzpicture}
%\caption{Strings over $\Sigma = \{a,b\}$ ordered by the containment relation (non-exhaustive)}
%\label{fig.string.contain}
%\end{figure}
%


\section{Logical Transduction}

Using model theory \citep{strother2019diss,yolyan2025logical,jardine2017local}, 
we model a phonological transformation from an underlying 
representation (UR) to a surface representation (SR) as a 
relation between input and output structures $\mathcal{A}_i \rightarrow \mathcal{A}_o$. 
There have been work adopting the same fashion on 
string-to-string transduction \refs, but the essentially 
the logic is not restricted to strings and can operate 
over any model-theoretic representations.
In this framework, we use a set of logic formula to derive
the output signature with reference to the input signature. 
Recall from the previous section that we defined an ASG 
signature as $\mathcal{A} = \langle D, \mathcal{T}, \alpha, \lhd \rangle$.
Given an input model with this signature, the output 
signature can be defined in terms of the input structure.
\begin{table}[ht!]
\centering
\footnotesize
\begin{tabular}{lcll}
\toprule
\textbf{Output} & & \textbf{Input} & \textbf{Interpretation} \\
\midrule

$\phi_{\textsf{Domain}}$
&
$:=$
&
$\textsc{True}$
&
Preserve all input elements
\\[0.5em]


\textsc{Copy}
&
$:=$
&
1&
Define Output domain the same size as Input
\\[0.5em]


$\phi_{\textsf{license}}(x)$
&
$:=$
&
$\exists y\, \alpha(x,y)$
&
Preserve only associated elements
\\[0.5em]

$T'(x)$
&
$:=$
&
$T(x)$
&
Preserve element labels
\\[0.5em]

$\lhd'(x,y)$
&
$\approx$
&
$\lhd(x,y)$
&
Preserve linear order within-tier
\\[0.5em]

$\alpha'(x,y)$
&
$:=$
&
$\bigl(
\alpha(x,y)\land
\neg \textsf{deassoc}(x,y)
\bigr)
\lor
\textsf{reassoc}(x,y)$
&
Update association relations
\\

\bottomrule
\end{tabular}

\caption{}
\label{tab.transduction}
\end{table}

As shown in Table~\ref{tab.transduction}, the transduction 
defines an output model with the same domain as the input 
structure, while filtering out unassociated elements. All 
unary predicates \(T(x)\) in the input signature are 
preserved in the corresponding output predicates \(T'(x)\), 
meaning that an element \(x\) has property \(T\) in the 
output iff it has property \(T\) in the input. Likewise, 
the within-tier order relation is preserved. The only 
change is the association relation \(\alpha\), which is defined such 
that an association line holds in the output iff either 
(i) it is present in the input and not removed by \textsf{deassoc}, or 
(ii) it is introduced by \textsf{reassoc}. 


\[
\begin{array}{rclclcl}
\alpha'(x,y)
&:=&
\bigl(
\alpha(x,y)
&\land&
\neg \textsf{deassoc}(x,y)
\bigr)
&\lor
&
\textsf{reassoc}(x,y)
\\[0.5em]

\text{output}
&=&
\text{unchanged}
&-&
\text{delinked}
&+&
\text{newly added}
\end{array}
\]

In this way, the problem of learning association-line 
transformations can be reduced to learning the logical 
formulas for \textsf{deassoc} and \textsf{reassoc}. 
Since these two predicates are combined through disjunction, 
two parallel learning processes take place simultaneously, 
and their outputs are combined to define the final 
association relation $\alpha'$.

We demonstrate this process using a common phenomenon in tone languages
known as \textit{contour formation}. When two vowels become adjacent,
the first vowel ($V_1$) may be deleted. Since the tone originally associated
with $V_1$ cannot remain unattached, the deletion triggers delinking of
$T_1$ from $V_1$ and reassociation of $T_1$ to the following vowel $V_2$.
As a result, the surviving vowel bears both tones in their original order,
which together surface as a contour tone.

From a logical perspective, this transformation can be characterized
through the interaction of the \textsf{deassoc},
\textsf{reassoc}, and $\phi_{\textsf{license}}$ functions.
The \textsf{deassoc} operation removes the association between the
deleted vowel and its tone, indicated by the crossed-out line in
Figure~\ref{fig:hiatus}. The \textsf{reassoc} operation then links the
tone of the deleted vowel ($T_1$) to the surviving vowel ($V_2$),
indicated by the dotted line in Figure~\ref{fig:hiatus}. Consequently,
$T_1$ and $V_1$ are no longer associated and therefore are not
licensed by $\phi_{\textsf{license}}$ (shown in gray in
Figure~\ref{fig:hiatus}). Thus, $T_1, T_2$ and $V_2$ remain in the
surface representation. If we define the relevant context as
$\textsf{hiatus}(x,y) := V(x) \land V(y) \land \lhd(x,y)$,
where $x$ is the first vowel ($V_1$) and $y$ is the second vowel ($V_2$),
then the transformation can be expressed as follows:
\[
\begin{array}{ll}
	\textsf{deassoc} = \alpha(t,v) \land \textsf{delete\_V (v)}\\
	
	\textsf{reassoc} = 	\textsf{survive\_V} \land\exists v\,(\alpha(t,v)\land \textsf{delete\_V})
\end{array}
\]

\begin{figure}[ht!]
	\centering
	
	\begin{minipage}[t]{0.48\textwidth}
		\centering
		
		\begin{tabular}{c c c}
			
			\textbf{UR} & $\Rightarrow$ & \textbf{SR} \\[0.5em]
			
			\begin{tikzpicture}[scale=0.7, every node/.style={inner sep=2pt}]
				\node [fill=gray!20](v1) at (0,0) {V$_1$};
				\node (v2) at (1.4,0) {V$_2$};
				
				\node (t1) at (0,1.3) {T$_1$};
				\node (t2) at (1.4,1.3) {T$_2$};
				
				\draw (v1)--(t1);
				\draw (v2)--(t2);
			\end{tikzpicture}
			
			&&
			
			\begin{tikzpicture}[scale=0.7, every node/.style={inner sep=2pt}]
				\node [fill=gray!20](v1) at (0,0) {V$_1$};
				\node (v2) at (1.4,0) {V$_2$};
				
				\node (t1) at (0,1.3) {T$_1$};
				\node (t2) at (1.4,1.3) {T$_2$};
				
				\draw (v1)--(t1) node[midway] {=};
				\draw (v2)--(t2);
				
				\draw[dotted, thick] (v2)--(t1);
			\end{tikzpicture}
			
		\end{tabular}
		
	\end{minipage}
	\hfill
	\begin{minipage}[t]{0.48\textwidth}
		\centering
		
		\begin{tabular}{c c c}
			
			\textbf{UR} & $\Rightarrow$ & \textbf{SR} \\[0.5em]
			
			$
			\begin{array}{c|cc}
				& T_1 & T_2\\
				\hline
				V_1 & \cellcolor{gray!75}1 & 0 \\
				V_2 & \cellcolor{gray!15}0 & 1 \\
			\end{array}
			$
			
			&&
			
			$
			\begin{array}{c|cc}
				& T_1 & T_2\\
				\hline
				V_1 & \cellcolor{gray!15}0 & 0 \\
				V_2 & \cellcolor{gray!75}1 & 1 \\
			\end{array}
			$
			
		\end{tabular}
		
	\end{minipage}
	
	\caption{
%		Vowel hiatus resolution and the corresponding adjacency matrices.
	}
	\label{fig:hiatus}
	
\end{figure}

The \textsf{deassoc} and \textsf{reassoc} operations are
equivalently reflected in the adjacency matrices.
Figure~\refs shows the input and output
matrices corresponding to the structures in
Figure~\ref{fig:hiatus}. As can be seen, the value
corresponding to the association between $T_1$ and $V_1$
changes from $1$ in the input to $0$ in the output,
reflecting delinking. Conversely, the value corresponding
to $T_1$ and $V_2$ changes from $0$ to $1$,
reflecting reassociation. In this way, the autosegmental
graph, the adjacency matrix, and the logical expressions
provide equivalent description of the same
transformation. Formally, \textsf{deassoc} and \textsf{reassoc}
can also be identified by  $ \textsf{deassoc}
= \{(i,j) \mid \mathrm{UR}_{ij}=1 \land \mathrm{SR}_{ij}=0\}$
and $\textsf{reassoc} =
\{(i,j) \mid \mathrm{UR}_{ij}=0 \land \mathrm{SR}_{ij}=1\}$.

\section{Association Learning}
 This section follows the framework of using logical transductions
to learn phonological mappings in \citet{yolyan2025logical},
which leverages the entailment relation
in a partially ordered space and updates hypotheses 
about the smallest contextual environments that trigger 
a featural change (a 
string-based input-output mapping that can be written as
$A \rightarrow B \,/\, C \_ D$).

Several aspects of this framework are particularly important for the present
work. First, phonological processes are represented as logical formulas
over model-theoretic representations. Second, the learning procedure relies on
grammatical entailment relations within a partially ordered hypothesis space,
with more general structures at the bottom and more specific structures higher
up \citep{chandleeBufia2019,rawski2021structure}. If an environment (which itself 
is a structure, thereafter we call it environment) is sufficient to trigger a
transformation, then any superfactor containing it also entails that
transformation. On the other hand, if an environment fails to trigger the
change, then all of its subfactors are likewise insufficient to trigger the
change. In other words, given a hypothesized environment, all data satisfying
this condition are expected to undergo the transformation (positive evidence),
while all data not satisfying this condition are expected not to undergo the
transformation (negative evidence). If some positive data fail to undergo the
transformation, then the hypothesis is too restrictive. Conversely, if some
negative data nevertheless undergo the transformation, then the hypothesis is
too general. The goal of learning is therefore to identify the smallest common 
environment that are contained in all positive evidence, but absent in all negative 
evidence. For example, post-nasal voicing can be written as 
\([-\mathrm{voi}] \rightarrow [+ \mathrm{voi}] \,/\, [+\mathrm{nasal}, -\mathrm{cont}] \_\_\),
meaning that whenever a segment specified as \([-\mathrm{voi}]\) occurs after a 
\([+\mathrm{nasal}, -\mathrm{cont}]\) segment, its voicing feature changes to 
\([+\mathrm{voi}]\). 

In autosegmental phonology, however, such contextual 
environments are not expressed linearly as \(A \rightarrow B \,/\, C \_ D\), 
but instead as structural relations over graphs. Nevertheless, because the 
learning framework applies generally to model-theoretic representations, it 
can naturally be extended to graph-based representations as well. 

As discussed previously, the learning problem reduces to identifying two types 
of structural changes: deassociation, where a matrix entry changes from 
\(1 \rightarrow 0\), and reassociation, where a matrix entry changes from 
\(0 \rightarrow 1\). The corresponding learning goal is therefore to identify 
the minimal environment around the changing cell that is both sufficient 
and exclusive in triggering the change observed in the positive evidence.


\[
0 \rightarrow 1 \;/\;
\begin{array}{ccc}
	? & ? & ? \\
	? & {\_} & ? \\
	? & ? & ?
\end{array}
\quad \text{or} \quad 
1 \rightarrow 0 \;/\;
\begin{array}{ccc}
	? & ? & ? \\
	? & {\_} & ? \\
	? & ? & ?
\end{array}
\]
ddTo learn ASGs, we adapt the bottom-up factor inference framework of
\citet{chandleeBufia2019} and the logical learning framework of
\citet{yolyan2025logical} to graph-based autosegmental representations
\refs. At the beginning of learning, the learner first
identifies changing association sites by comparing the input and output
association matrices. If an association value changes from
\(1 \rightarrow 0\), the change is classified as deassociation; if it
changes from \(0 \rightarrow 1\), it is classified as reassociation.
The learner then initializes the search with the corresponding most
general logical hypothesis: $\neg \alpha(x,y)$ for reassociation and
$\exists x \exists y \alpha(x,y)$ for deassociation \footnote{We treat deassociation and
	reassociation as two separate learning problems and combine the
	resulting hypotheses through disjunction.}.

\begin{algorithm}[th!]
	\small
	\LinesNumbered
	
	\KwData{
		$
		D	=	\{(I_1,O_1),\dots,(I_n,O_n)\},
		k, V^+ = \{\, I_n \mid I_n \neq O_n \,\},
		V^- = \{\, I_n \mid I_n = O_n \,\} $
	}
%	
%	\KwOut{
%		minimal structural environment $c$
%	}
	
	$P \gets \emptyset$\;
%	\tcp*{pruned set}

	$Q \gets \{\alpha(x,y)\}$ or 	$Q \gets \{\neg \alpha(x,y)\}$\; 	

	\While{$Q \neq \emptyset$}{
		
		$c \gets Q.\mathrm{Dequeue}()$\;
	
		\If{$\forall g \in V^+,\; c \sqsubseteq g$}{
			
			\If{$\forall g \in V^-,\; c \not\sqsubseteq g$}{
				
				\Return{$c$}\;
			}
			
			\Else{
				
				\ForEach{$s \in \textsc{NextSupFact}(c) \land |s| < k \land \forall p \in P,\; s \not\sqsubseteq p$}{
					
				
						
						$Q.\mathrm{Enqueue}(s)$\;
					
				}
			}
		}
		
		\Else{
		
			$P \gets P \cup \{c\}$\;
		}
	}
	\Return{$\bot$}\;
	
	\caption{Learning Association Changes}
\end{algorithm}


% explain the algo ! fix Q? 
The input data consists of a finite set of input-output ASG pairs,
which can be partitioned into positive and negative evidence, denoted
$V^+$ and $V^-$ respectively. It also requires a bound $k$ that limits
the graph size. 

The algorithm conducts a breadth-first search over a partially ordered
hypothesis space. At each step, a candidate hypothesis $c$ is tested 
against the positive evidence: if $c$ is not contained in
all positive examples, then the current hypothesis is incompatible with the
observed transformation and thus is pruned. Otherwise, the learner checks
 $c$ against all negative examples: if none of data contains c, then
 we return the current structure, otheewise,  it
is considered too general. The learner therefore generates more
specific superfactors using \textsc{NextSupFact}.
Because each application of \textsc{NextSupFact} strictly increases the
size of the hypothesis and the search is bounded by $k$, the hypothesis
space is finite and the algorithm terminates. Moreover, since
hypotheses are explored breadth-first, the first valid hypothesis
returned by the learner is guaranteed to be the minimal structural
environment consistent with all positive and negative evidence.
One crucial component of the learning procedure is
\textsc{NextSupFact}. There are many possible ways to generate
superfactors while still covering the same hypothesis space. Here, we
propose one possible method that explicitly connects graph-theoretic
representations with their logical counterparts.

From a graph-theoretic perspective, the default first step is to expand
the domain by introducing neighboring nodes. Such neighboring nodes may
be added on either tier in any successor relation to existing
nodes. Table~\ref{tab:nextsupfact} presents a non-exhaustive set of examples
illustrating the generation of \textsc{NextSupFact} hypotheses starting
from an initial deassociation hypothesis. By adding neighbor, one can 
expect a possible candidate with an additional node \(z\) following \(x\). 
Other possible configurations include \(x\) following \(z\),
\(z\) following \(x\), \(z\) following \(y\), or \(y\) following \(z\).
Once additional nodes are introduced and valid association docking sites
become available, new association lines may be added between nodes.
Finally, the learner may further specialize the hypothesis by specifying
node labels or feature values~\footnote{For a more detailed description of \textsc{NextSupFact} operations over ASGs, see \refs.}.



\begin{table}[ht!]
	\centering
	\small
	\begin{tabular}{p{2.8cm}p{2.5cm}p{2.5cm}p{6cm}} 
		\toprule
		\textbf{Operation} & \textbf{Matrix} & \textbf{Graph} & \textbf{Logic} \\
		\midrule
		Initial hypothesis 
		&
		$
		\begin{array}{c|c}
			& x  \\
			\hline
			y & 1 
		\end{array}
		$
		&
		\begin{tikzpicture}[baseline=(current bounding box.center)]
			\draw[fill=black] (0,0) circle (2pt);
			\draw[fill=black] (0,1) circle (2pt);
			
			\draw (0,0) -- (0,1);
			
			\node[right] at (0,0) {$x$};
			\node[right] at (0,1) {$y$};
		\end{tikzpicture}
		&
		$\alpha(x,y)$
		\\[2em]
		
		Add neighbor
		&
		$\begin{array}{c|cc}
			& x & z \\
			\hline
			y & 1 & 
		\end{array}$		&
		\begin{tikzpicture}[baseline=(current bounding box.center)]
			\node (x) at (0,0) {};
			\node (y) at (0,1) {};
			\node (z) at (1,0) {};
			
			\draw[fill=black] (x) circle (2pt);
			\draw[fill=black] (y) circle (2pt);
			\draw[fill=black] (z) circle (2pt);
			
			\draw (x) -- (y);
			\draw[dashed] (x) -- (z);
			
			\node[right] at (x) {$x$};
			\node[right] at (y) {$y$};
			\node[right] at (z) {$z$};
		\end{tikzpicture}
		&
		$
		\alpha(x,y) \land \exists z \lhd(x,z)$
		\\[2em]
		
		Add association
		&
		$
		\begin{array}{c|cc}
			& x & z \\
			\hline
			y & 1 & 1
		\end{array}
		$
		&
		\begin{tikzpicture}[baseline=(current bounding box.center)]
			\node (x) at (0,0) {};
			\node (y) at (0,1) {};
			\node (z) at (1,0) {};
			
			\draw[fill=black] (x) circle (2pt);
			\draw[fill=black] (y) circle (2pt);
			\draw[fill=black] (z) circle (2pt);
			
			\draw (x) -- (y);
			\draw (x) -- (z);
			\draw[dashed] (y) -- (z);			
			\node[right] at (x) {$x$};
			\node[right] at (y) {$y$};
			\node[right] at (z) {$z$};
		\end{tikzpicture}
		&
		$		\alpha(x,y) \land \exists z\,\big(\lhd(x,z) \land \alpha(y,z)\, \big)$
		\\[2em]
		Add label
		&
		$
		\begin{array}{c|cc}
			& H & z \\
			\hline
			y & 1 & 1
		\end{array}
		$
		&
		\begin{tikzpicture}[baseline=(current bounding box.center)]
			\node (x) at (0,0) {};
			\node (y) at (0,1) {};
			\node (z) at (1,0) {};
			
			\draw[fill=black] (x) circle (2pt);
			\draw[fill=black] (y) circle (2pt);
			\draw[fill=black] (z) circle (2pt);
			
			\draw (x) -- (y);
			\draw (x) -- (z);
			\draw[] (y) -- (z);			
			\node[left] at (x) {H};
			\node[right] at (y) {$y$};
			\node[right] at (z) {$z$};
		\end{tikzpicture}
		&
		$		\alpha(x,y) \land \exists z\,\big(\lhd(x,z) \land \alpha(y,z)\, \land H(x)\big)$
		\\[2em]
		\bottomrule
	\end{tabular}
	
	\caption{Examples of operations used in generating 
		\textsc{NextSupFact} hypotheses across matrix, graph, and logical 
		representations. }
	\label{tab:nextsupfact}
\end{table}

As can be seen, this provides a clear correspondence between
graph-theoretic and logical representations. New quantifiers are
introduced through successor relations, corresponding to the addition of
neighboring nodes in the graph structure. Within the same quantifier
scope, the learner may further specialize the hypothesis through
conjunction, including adding association relations (corresponding to
changing matrix cells) or specifying node labels and feature values.

\section{Case Study}
\subsection{Contour Formation}

\section{Conclusion} This paper has argued that vowel-deletion–induced 
tonal processes can be uniformly captured as logical transductions 
over autosegmental representations (ARs). The observed typological 
variation follows from modifying a single relation, $\alpha$, through 
the interaction of \textsf{deassoc} and \textsf{reassoc}. 

A limitation of the current analysis is that the transduction relies 
on quantification to capture across-tier dependencies. An important 
direction for future work is to determine whether these processes can 
be simplified in a quantifier-free (QF) process.
This work also extends naturally to learning over ARs. Recent work 
by \citet{yolyan2025logical} demonstrates that logical transductions 
are learnable by using a bottom-up factorial inference algorithms\citep{chandleeBufia2019}. This opens up more discussions about how linguistically meaningful 
representations can facilitate grammar inference.


\bibliographystyle{apalike}
\bibliography{ref}

\end{document}